% Mathematical TeX extracted from the cited web pages, reformatted by the assistant. Not the native full chapter.
\begin{lemma}[10.51.2, Artin--Rees]
Suppose that $R$ is Noetherian, $I\subset R$ an ideal. Let $N\subset M$ be finite $R$-modules. There exists a constant $c>0$ such that $I^nM\cap N=I^{n-c}(I^cM\cap N)$ for all $n\geq c$.
\end{lemma}
\begin{proof}
Consider the ring $S=R\oplus I\oplus I^2\oplus\ldots=\bigoplus_{n\geq0}I^n$. Convention: $I^0=R$. Multiplication maps $I^n\times I^m$ into $I^{n+m}$ by multiplication in $R$. Note that if $I=(f_1,\ldots,f_t)$ then $S$ is a quotient of the Noetherian ring $R[X_1,\ldots,X_t]$. The map just sends the monomial $X_1^{e_1}\ldots X_t^{e_t}$ to $f_1^{e_1}\ldots f_t^{e_t}$. Thus $S$ is Noetherian.

Similarly, consider the module $M\oplus IM\oplus I^2M\oplus\ldots=\bigoplus_{n\geq0}I^nM$. This is a finitely generated $S$-module. Namely, if $x_1,\ldots,x_r$ generate $M$ over $R$, then they also generate $\bigoplus_{n\geq0}I^nM$ over $S$. Next, consider the submodule $\bigoplus_{n\geq0}I^nM\cap N$. This is an $S$-submodule, as is easily verified. By Lemma 10.51.1 it is finitely generated as an $S$-module, say by $\xi_j\in\bigoplus_{n\geq0}I^nM\cap N$, $j=1,\ldots,s$. We may assume by decomposing each $\xi_j$ into its homogeneous pieces that each $\xi_j\in I^{d_j}M\cap N$ for some $d_j$. Set $c=\max\{d_j\}$. Then for all $n\geq c$ every element in $I^nM\cap N$ is of the form $\sum h_j\xi_j$ with $h_j\in I^{n-d_j}$. The lemma now follows from this and the trivial observation that $I^{n-d_j}(I^{d_j}M\cap N)\subset I^{n-c}(I^cM\cap N)$.
\end{proof}

\begin{lemma}[10.97.1]
Let $I$ be an ideal of a Noetherian ring $R$. Denote ${}^\wedge$ completion with respect to $I$.
\begin{enumerate}
\item If $K\to N$ is an injective map of finite $R$-modules, then the map on completions $K^\wedge\to N^\wedge$ is injective.
\item If $0\to K\to N\to M\to0$ is a short exact sequence of finite $R$-modules, then $0\to K^\wedge\to N^\wedge\to M^\wedge\to0$ is a short exact sequence.
\item If $M$ is a finite $R$-module, then $M^\wedge=M\otimes_RR^\wedge$.
\end{enumerate}
\end{lemma}
\begin{proof}
Setting $M=N/K$ we find that part (1) follows from part (2). Let $0\to K\to N\to M\to0$ be as in (2). For each $n$ we get the short exact sequence
\[
0\to K/(I^nN\cap K)\to N/I^nN\to M/I^nM\to0.
\]
By Lemma 10.87.1 we obtain the exact sequence
\[
0\to\varprojlim K/(I^nN\cap K)\to N^\wedge\to M^\wedge\to0.
\]
By the Artin--Rees Lemma 10.51.2 we may choose $c$ such that $I^nK\subset I^nN\cap K\subset I^{n-c}K$ for $n\geq c$. Hence $K^\wedge=\varprojlim K/I^nK=\varprojlim K/(I^nN\cap K)$ and we conclude that (2) is true.

Let $M$ be as in (3) and let $0\to K\to R^{\oplus t}\to M\to0$ be a presentation of $M$. We get a commutative diagram
\[
\xymatrix@C=10pt{
&K\otimes_RR^\wedge\ar[r]\ar[d]&R^{\oplus t}\otimes_RR^\wedge\ar[r]\ar[d]&M\otimes_RR^\wedge\ar[r]\ar[d]&0\\
0\ar[r]&K^\wedge\ar[r]&(R^{\oplus t})^\wedge\ar[r]&M^\wedge\ar[r]&0
}
\]
The top row is exact, see Section 10.39. The bottom row is exact by part (2). By Lemma 10.96.1 the vertical arrows are surjective. The middle vertical arrow is an isomorphism. We conclude (3) holds by the Snake Lemma 10.4.1.
\end{proof}

\begin{lemma}[10.97.2]
Let $I$ be an ideal of a Noetherian ring $R$. Denote ${}^\wedge$ completion with respect to $I$.
\begin{enumerate}
\item The ring map $R\to R^\wedge$ is flat.
\item The functor $M\mapsto M^\wedge$ is exact on the category of finitely generated $R$-modules.
\end{enumerate}
\end{lemma}
\begin{proof}
Consider $J\otimes_RR^\wedge\to R\otimes_RR^\wedge=R^\wedge$ where $J$ is an arbitrary ideal of $R$. According to Lemma 10.97.1 this is identified with $J^\wedge\to R^\wedge$ and $J^\wedge\to R^\wedge$ is injective. Part (1) follows from Lemma 10.39.5. Part (2) is a reformulation of Lemma 10.97.1 part (2).
\end{proof}

\begin{lemma}[10.97.3]
Let $I$ be an ideal of a Noetherian ring $R$. Denote $R^\wedge$ the completion of $R$ with respect to $I$. If $I$ is contained in the Jacobson radical of $R$, then the ring map $R\to R^\wedge$ is faithfully flat. In particular, if $(R,\mathfrak m)$ is a Noetherian local ring, then the completion $\varprojlim_nR/\mathfrak m^n$ is faithfully flat.
\end{lemma}
\begin{proof}
By Lemma 10.97.2 it is flat. The composition $R\to R^\wedge\to R/I$ where the last map is the projection map $R^\wedge\to R/I$ shows that any maximal ideal of $R$ is in the image of $\Spec(R^\wedge)\to\Spec(R)$. Hence the map is faithfully flat by Lemma 10.39.15.
\end{proof}

